\section{Chương 3: Ứng dụng thực tế}
\label{sec:applications}

\begin{table*}[!t]
\centering\small
\caption{Mười luật kết hợp đứng đầu theo lift trên Groceries. Số giỏ tính cho cả hai vế. Bảng hiển thị các chỉ số đã làm tròn.}
\label{tab:groceries-top10-m2}
\TableGroceries
\end{table*}

\subsection{Dữ liệu và cách sinh luật trên Groceries}
Groceries gồm 9\,835 giỏ và 169 nhóm sản phẩm. Nhóm dùng một phiên
bản cố định của tệp CSV trong kho \emph{Machine Learning with
R}~\cite{groceries_source_b}. Tài liệu \texttt{arules} mô tả dữ liệu
được thu trong 30 ngày~\cite{groceries_documentation_b}. Mỗi dòng
tệp nguồn là một giỏ. Chương trình gán mã giỏ và chuyển tệp thành
43\,367 cặp \texttt{basket\_id,item\_name}. Julia chuẩn hóa tên về
dạng NFC của Unicode, chuyển sang chữ thường và chuẩn hóa khoảng
trắng, sau đó gán mã sản phẩm theo thứ tự từ điển. Mỗi sản phẩm chỉ
tính một lần trong một giỏ. Hai giỏ giống nội dung vẫn được đếm riêng.

Nhóm chọn minsup $0{,}005$ để giữ các tổ hợp ít gặp nhưng vẫn xuất
hiện trong ít nhất $\lceil0{,}005\times9835\rceil=50$ giỏ. Ngưỡng
độ tin cậy tối thiểu (minimum confidence, minconf) là $0{,}60$,
tức hậu kiện phải xuất hiện trong ít nhất 60\% số giỏ thỏa tiền
kiện. Hai bản ECLAT cùng tìm được 1\,001 tập phổ biến.

Gọi $A$ là tiền kiện và $B$ là hậu kiện, cả hai khác rỗng và rời
nhau. Ký hiệu $n$ là số giỏ, còn $\sigma(X)$ là số giỏ chứa $X$.
Độ hỗ trợ của luật là tỷ lệ giỏ chứa cả hai vế. Độ tin cậy là tỷ lệ
giỏ chứa $B$ trong số những giỏ đã chứa $A$. Lift so độ tin cậy
với tỷ lệ xuất hiện của $B$ trong toàn bộ dữ liệu.
\begin{align}
\operatorname{supp}(A\Rightarrow B)&=\sigma(A\cup B)/n,\\
\operatorname{conf}(A\Rightarrow B)&=\sigma(A\cup B)/\sigma(A),\\
\operatorname{lift}(A\Rightarrow B)&=\operatorname{conf}(A\Rightarrow B)/(\sigma(B)/n).
\end{align}
Từ mỗi tập phổ biến, chương trình xét mọi cách chia thành hai vế
và dùng tỷ số chính xác của các số đếm để lọc theo độ tin cậy.
Có 22 luật đạt ngưỡng, xếp giảm dần theo lift, độ tin cậy và độ
hỗ trợ trước khi làm tròn. Nếu cả ba chỉ số bằng nhau, chương
trình ưu tiên hợp hai vế ngắn hơn rồi xét thứ tự từ điển của
mã ở hai vế.

\subsection{Ý nghĩa kinh doanh của các luật trên Groceries}
Luật đứng đầu trong Bảng~\ref{tab:groceries-top10-m2} có tiền kiện
gồm \emph{citrus fruit}, \emph{root vegetables} và \emph{whole milk}.
Trong 90 giỏ có đủ ba nhóm hàng này, 57 giỏ cũng có \emph{other
vegetables}. Độ tin cậy 63,33\% cao hơn tỷ lệ nền
$1903/9835=19{,}35\%$, tương ứng lift 3,273. Nhóm rau này là một
ứng viên để thử gợi ý mua kèm cho các giỏ thỏa tiền kiện. Độ hỗ
trợ 0,5796\%, tức 57 giỏ trên toàn bộ dữ liệu, giới hạn phạm vi
mà số liệu trực tiếp mô tả.

Luật hạng 5 gợi ý \emph{whole milk} khi giỏ chứa \emph{root
vegetables}, \emph{tropical fruit} và \emph{yogurt}. Độ tin cậy
là $56/80=70\%$ so với tỷ lệ nền của sữa là 25,55\%. Luật có độ
hỗ trợ 0,5694\% và lift 2,740. Với cùng hậu kiện, luật hạng 9
chỉ cần \emph{butter} và \emph{yogurt}. Có $92/144=63{,}89\%$
giỏ thỏa hậu kiện, với lift 2,500 và độ hỗ trợ 0,9354\%.
Tiền kiện ngắn và độ hỗ trợ cao hơn luật hạng 1 khiến luật
hạng 9 là một lựa chọn khi ưu tiên phạm vi thử nghiệm.

Trong 9\,835 giỏ, có 580 giỏ thỏa tiền kiện của ít nhất một
trong mười luật đứng đầu. Có 353 giỏ thỏa cả hai vế của ít nhất
một luật, chiếm 3,59\% tổng số giỏ. Vì một giỏ có thể thỏa nhiều
luật, hai con số này đếm số giỏ phân biệt để mô tả phạm vi áp
dụng của các gợi ý.

Lift đo mức đồng xuất hiện so với giả thiết độc lập. Độ hỗ trợ
thấp và việc chọn luật trên cùng dữ liệu có thể làm thứ hạng
thiếu ổn định. Dữ liệu thiếu giá, thời điểm và nhóm đối chứng,
nên các luật phù hợp để đặt giả thuyết bán chéo. Một thử nghiệm
A/B có thể phân ngẫu nhiên các giỏ đủ điều kiện vào nhóm nhận
gợi ý và nhóm đối chứng rồi so tỷ lệ thêm sản phẩm giữa hai nhóm.

\subsection{Mã hóa và chuẩn bị dữ liệu Splice DNA}
\begin{table}[!tb]
\centering
\caption{Số mẫu trước và sau khi lọc theo lớp. Ngưỡng đếm là 20\% số mẫu giữ lại, làm tròn lên.}
\label{tab:splice-preparation-b}
\begin{tabular}{lrrrr}
\toprule
Lớp & Ban đầu & Giữ lại & Loại & Ngưỡng đếm \\
\midrule
EI & 767 & 765 & 2 & 153 \\
IE & 768 & 766 & 2 & 154 \\
N & 1\,655 & 1\,654 & 1 & 331 \\
\bottomrule
\end{tabular}
\end{table}

Bộ Splice của UCI~\cite{splice} có 3\,190 chuỗi, mỗi chuỗi gồm
60 vị trí và có mã quan sát cùng nhãn. EI là ranh giới từ exon
sang intron, IE là ranh giới từ intron sang exon, còn N chỉ
các chuỗi ngoài hai lớp điểm nối. Nhóm dùng bản lưu tại Đại
học Szeged~\cite{splice_mirror_b}, xác minh SHA-256 và giữ
mã quan sát theo thứ tự dòng. Tệp \texttt{splice.data} trong
ZIP chính thức của UCI có cùng giá trị băm và là nguồn dự phòng.

Cửa sổ gồm 12 nucleotide tại chỉ số \texttt{25:36} trong Julia,
tương ứng $-6,\ldots,-1,+1,\ldots,+6$. Với EI và IE, tọa độ
tính từ ranh giới đã chú thích. Với N, tọa độ tính quanh tâm
cửa sổ đối chứng. Mỗi mục là một cặp vị trí và nucleotide. Bốn
khả năng A, C, G và T tại mỗi vị trí tạo thành 48 mã cố định.
Chẳng hạn, cùng một nucleotide ở hai phía ranh giới sẽ có hai
mã khác nhau. Nhãn lớp được lưu riêng khỏi giao dịch.

Mỗi chuỗi tạo một giao dịch nếu cả 12 vị trí đều xác định là
A, C, G hoặc T. Ký hiệu mơ hồ ngoài cửa sổ vẫn được chấp nhận.
Bảng~\ref{tab:splice-preparation-b} cho thấy năm mẫu bị loại,
còn lại 3\,185 mẫu.
Nhóm khai thác riêng từng lớp để xác định các tổ hợp phổ biến
trong từng loại ranh giới và trong nhóm đối chứng. Minsup 20\%
áp cùng yêu cầu về tỷ lệ xuất hiện dù số mẫu giữa các lớp khác
nhau. Chương trình đổi tỷ lệ này thành ngưỡng đếm theo số mẫu
được giữ ở mỗi lớp.

\subsection{So sánh các mẫu nucleotide giữa ba lớp}
\begin{table}[!tb]
\centering\small
\caption{Độ hỗ trợ (\%) của bốn mẫu trên EI, IE và N. Hạng tính trong danh sách 463 ứng viên.}
\label{tab:dna-m2}
\TableSplice
\end{table}

Hai bản ECLAT khớp kết quả trên từng lớp, với 381 tập ở EI,
209 tập ở IE và 48 tập ở N. Tập dài nhất ở EI có 6 mục. Để so
sánh, nhóm lấy hợp các tập phổ biến của ba lớp và
giữ những tập dài từ 2 đến 4 mục, thu được 463 ứng viên. Sau đó,
chương trình đếm lại mỗi ứng viên trên cả ba lớp. Trong 864
trường hợp, mẫu có xuất hiện ở một lớp nhưng số lần xuất hiện
chưa đạt ngưỡng của lớp đó. Các số đếm này vẫn góp vào chênh
lệch độ hỗ trợ.

Nhóm xếp hạng theo chênh lệch độ hỗ trợ giảm dần,
$\max(s_{EI},s_{IE},s_N)-\min(s_{EI},s_{IE},s_N)$, trong đó
$s_c$ là tỷ lệ mẫu giữ lại ở lớp $c$ có chứa ứng viên. Khi
chênh lệch bằng nhau, chương trình xét độ hỗ trợ lớn nhất theo
thứ tự giảm dần, sau đó đến độ dài và mã mục theo thứ tự tăng
dần. Bảng~\ref{tab:dna-m2} trích bốn mẫu đại diện cho độ dài
2, 3 và 4 mục, với thứ hạng trong danh sách đầy đủ.

Cặp $(+1,G),(+2,T)$ xuất hiện trong $758/765=99{,}08\%$ mẫu
EI, $108/766=14{,}10\%$ mẫu IE và $98/1654=5{,}93\%$ mẫu N.
Ở phía bên kia ranh giới, cặp $(-2,A),(-1,G)$ xuất hiện trong
$762/766=99{,}48\%$ mẫu IE, $388/765=50{,}72\%$ mẫu EI và
$127/1654=7{,}68\%$ mẫu N. Tính từ tỷ lệ trước khi làm tròn,
chênh lệch giữa EI và N đối với cặp đầu là 93,16 điểm phần
trăm, còn chênh lệch giữa IE và N đối với cặp sau là 91,80
điểm phần trăm.

GT ngay sau ranh giới EI và AG ngay trước ranh giới IE phù
hợp với hai cặp nucleotide thường gặp ở vị trí cho và vị trí
nhận của intron ở động vật có vú~\cite{burset2000}. Cả hai cặp
cũng xuất hiện trong lớp N, vì vậy sự có mặt của từng cặp
chưa đủ để nhận diện điểm nối. Mẫu 4 mục ở hạng 13 có số đếm
lần lượt là 514 ở EI, 22 ở IE và 2 ở N. Khi thêm điều kiện
về nucleotide ở các vị trí khác, độ hỗ trợ ở EI giảm còn
67,19\%.

Các tỷ lệ trên mô tả sự đồng xuất hiện trong bộ chuỗi đã căn
chỉnh. Cửa sổ, ngưỡng và cách xử lý ký hiệu mơ hồ ảnh hưởng
đến tập mẫu thu được. Để đánh giá khả năng nhận diện điểm nối
trên chuỗi mới, cần kiểm tra trên dữ liệu độc lập.

\subsection{Kiểm chứng và khả năng tái lập}
Chương trình Python đếm lại độ hỗ trợ trực tiếp từ các giao
dịch và đối chiếu thứ hạng. Với cùng ngưỡng đếm, SPMF 2.64b
khớp toàn bộ tập và độ hỗ trợ của 1\,001 tập trên Groceries
cùng 381, 209 và 48 tập trên ba lớp Splice. README dẫn đến
ba notebook có ghi mã phiên bản của mỗi lần chạy. Hai notebook
ứng dụng tải dữ liệu, chạy bản Julia của nhóm và tạo bảng, hình
từ kết quả mới.
